\section{Conditional State Posterior: Kalman Filter and FFBS}

\subsection{Unified Linear-Gaussian Form}
Conditioning on static parameters, each model block reduces to
\begin{align}
\bm x_t\given\bm x_{t-1} &\sim \N(\bm G_t\bm x_{t-1},\bm Q_t), \label{eq:lgssm_state}\\
y_{t,n}\given\bm x_t &\sim \N(\bm h_{t,n}^\T\bm x_t,r_{t,n}),\quad n=1,\dots,n_t. \label{eq:lgssm_obs}
\end{align}
Here $\bm x_t$ is whichever stacked state is active ($\bm\alpha_t$, $\bm\beta_t$, or block-stacked variants).

\subsection{Sequential Kalman Assimilation at Fixed Time}
Prediction:
\begin{align}
\bm a_t &= \bm G_t\bm m_{t-1},
& \bm R_t &= \bm G_t\bm C_{t-1}\bm G_t^\T + \bm Q_t.
\end{align}
For each scalar observation $(y_{t,n},\bm h_{t,n},r_{t,n})$:
\begin{align}
f_{t,n} &= \bm h_{t,n}^\T \bm R_{t,n-1}\bm h_{t,n}+r_{t,n}, \label{eq:kf_f}\\
\bm K_{t,n} &= \bm R_{t,n-1}\bm h_{t,n} f_{t,n}^{-1}, \label{eq:kf_K}\\
\bm m_{t,n} &= \bm m_{t,n-1}+\bm K_{t,n}(y_{t,n}-\bm h_{t,n}^\T\bm m_{t,n-1}), \label{eq:kf_m}\\
\bm C_{t,n} &= \bm C_{t,n-1}-\bm K_{t,n}\bm h_{t,n}^\T\bm C_{t,n-1}. \label{eq:kf_C}
\end{align}
Initialization at time $t$: $\bm m_{t,0}=\bm a_t$, $\bm C_{t,0}=\bm R_t$.
Set filtered moments $\bm m_t=\bm m_{t,n_t}$, $\bm C_t=\bm C_{t,n_t}$.
Thus, the second index $n$ denotes within-time assimilation stage, while
$(\bm m_t,\bm C_t)$ denote post-assimilation filtered moments at time $t$.

\subsection{FFBS Simulation Smoother}
Backward simulation starts at $\bm x_T\sim\N(\bm m_T,\bm C_T)$.
For $t=T-1,\dots,0$,
\begin{align}
\bm B_t &= \bm C_t\bm G_{t+1}^\T\bm R_{t+1}^{-1}, \label{eq:ffbs_B}\\
\bm h_t &= \bm m_t + \bm B_t(\bm x_{t+1}-\bm a_{t+1}), \label{eq:ffbs_h}\\
\bm H_t &= \bm C_t - \bm B_t\bm R_{t+1}\bm B_t^\T, \label{eq:ffbs_H}\\
\bm x_t\given\bm x_{t+1},\text{data},\Theta &\sim \N(\bm h_t,\bm H_t). \label{eq:ffbs_cond}
\end{align}

\subsection{Computability Constraints}
All updates use Cholesky solves of $f_{t,n}$ (scalar) or small SPD matrices.
No inversion of large $n_t\times n_t$ matrices is required when observations are assimilated sequentially.

\subsection{Ragged-horizon forecast FFBS (Model C)}
For Model C with source-specific horizons $K_j$, the lead-specific active set
$\mathcal A_k=\{j:k\le K_j\}$ induces reduced state dimension $p_k=q(1+|\mathcal A_k|)$.
Two equivalent implementations are valid:
\begin{enumerate}
\item True transdimensional recursion in $\R^{d_k^{(f)}}$ with lead-dependent
      $(\bm M_{T+k}^{(f)},\bm Q_{T+k}^{(f)},\bm h_{T+k,n}^{(f)})$, where
      $d_k^{(f)}=p_k$ in the reduced specification and $d_k^{(f)}=p_k+1+m$
      in the augmented specification.
\item Fixed-dimension embedding in $\R^{q(1+J_f)}$ (reduced) or
      $\R^{q(1+J_f)+1+m}$ (augmented) with deterministic selection/projection
      matrices that map the embedded state to the active state used by the
      forecast likelihood at each lead.
\end{enumerate}
The fixed-dimension form is algebraically equivalent to the transdimensional form
provided inactive coordinates are projected out of both the measurement operators and
the state innovation terms at each lead.
For Gibbs updates, $\bm Q_{T+k}^{(f)}$ is the sampled lead-specific covariance
$\bm W_{T+k}^{(f)}$.
For mean-field VB, the pseudo-transition at lead $T+k$ uses the precision
$\E_q[(\bm W_{T+k}^{(f)})^{-1}]$, or equivalently the derived pseudo-covariance
$\bar{\bm W}_{T+k}^{(f)}:=\E_q[(\bm W_{T+k}^{(f)})^{-1}]^{-1}$ in covariance-form
recursions; this pseudo-covariance is not $\E_q[\bm W_{T+k}^{(f)}]$.
